\chapter{Environmental First Aid}\label{ch:environmental}

\section{Heat-Related Illnesses}

Heat-related illnesses occur when the body cannot properly regulate temperature in hot environments, especially with high humidity and inadequate hydration.

\subsection{Heat Cramps}

Mild form characterized by painful muscle spasms during or after strenuous activity in heat.

\begin{itemize}[leftmargin=*]
  \item \textbf{Symptoms}: Painful muscle cramps (often legs, abdomen, arms), heavy sweating, normal body temperature
  \item \textbf{Management}: Rest in cool place; drink water or sports drinks containing electrolytes; gently stretch and massage affected muscles; avoid strenuous activity for several hours afterward
  \item Seek medical attention if cramps do not improve within an hour, person has heart condition or is on low-sodium diet, or vomiting develops
\end{itemize}

\subsection{Heat Exhaustion}

More serious than heat cramps but less severe than heat stroke.

\begin{itemize}[leftmargin=*]
  \item \textbf{Symptoms}: Heavy sweating, weakness, dizziness, nausea, headache, cool/moist skin, rapid but weak pulse, normal to slightly elevated body temperature (up to 104 degrees Fahrenheit / 40 degrees Celsius)
  \item \textbf{Management}: Move person to cooler environment; have them lie down and elevate legs slightly; remove excess clothing; apply cool wet cloths or mist with water while fanning; give small sips of cool water if conscious and able to swallow (do not force fluids); monitor closely
  \item If no improvement within 30 minutes, symptoms worsen, person vomits, or becomes unresponsive -- CALL EMERGENCY SERVICES IMMEDIATELY
  \item NOTE: Heat exhaustion CAN progress to heat stroke, which is a life-threatening emergency requiring rapid cooling and immediate medical care
\end{itemize}

\subsection{Heat Stroke (Medical Emergency)}

The most serious heat illness where the body's temperature regulation system fails.

\begin{itemize}[leftmargin=*]
  \item \textbf{Symptoms}: Body temperature greater than 104 degrees Fahrenheit (40 degrees Celsius), hot/dry skin OR profuse sweating (depending on type), altered mental state (confusion, agitation, slurred speech, combativeness, seizures, unconsciousness), rapid strong pulse, potential loss of consciousness
  \item \textbf{Management}: CALL EMERGENCY SERVICES IMMEDIATELY. Rapid cooling is critical: immerse person in cold water if possible (tub, pool, lake); otherwise pour/continuously spray copious amounts of cool water over skin and fan aggressively; place ice packs on neck, armpits, and groin (major blood vessels near surface). DO NOT give fluids to an unconscious person. Monitor breathing continuously and be prepared to perform CPR if necessary \cite{aha_arc_2024,cdc_heat}.
  \item \textbf{Important}: Do NOT wait for professional help to begin cooling. Every minute counts in preventing brain damage or death.
\end{itemize}

\textbf{WARNING:} \textbf{NEVER} leave a person with suspected heat stroke alone. Continuously monitor breathing and consciousness level. Even if they seem to recover, seek medical evaluation as there may be internal organ damage.

\noindent\textbf{Heat waves in India}: Severe heat waves, mainly from April to June, are a major public health risk in India, particularly for outdoor workers, the elderly, children, and people without access to cooling. The National Disaster Management Authority (NDMA) issues district-level heat wave warnings and action plans \cite{ndma}. During a heat wave, avoid outdoor activity between 12 noon and 3 p.m., drink water regularly, and check on vulnerable neighbours \cite{cdc_heat}. See Chapter~\ref{ch:disaster} for heat wave preparedness details.

\section{Cold-Related Illnesses}

\subsection{Hypothermia}

Occurs when the body loses heat faster than it can produce, causing dangerously low body temperature (below 95 degrees Fahrenheit / 35 degrees Celsius).

\subsubsection{Stages and Symptoms}

\begin{center}
\small
\begin{tabularx}{\textwidth}{@{}>{\raggedright\arraybackslash}l>{\raggedright\arraybackslash}X>{\raggedright\arraybackslash}X@{}}
\hline
\textbf{Stage} & \textbf{Temperature Range} & \textbf{Signs \& Symptoms} \\
\hline
Mild & 90-95 degrees Fahrenheit (32-35 degrees Celsius) & Shivering, numbness, clumsiness, slurred speech, fatigue, coordination problems \\
Moderate & 82-90 degrees Fahrenheit (28-32 degrees Celsius) & Intense shivering stops, confusion, drowsiness, slow/shallow breathing, weak pulse, poor coordination, apathy \\
Severe & Below 82 degrees Fahrenheit (28 degrees Celsius) & No shivering, very slow/irregular breathing and pulse, unresponsive/comatose, pupils dilated, risk of cardiac arrest \\
\hline
\end{tabularx}
\end{center}

\textbf{WARNING:} \textbf{CAUTION:} "Paradoxical Undressing": As hypothermia progresses, some people feel falsely warm and begin removing clothing -- this accelerates heat loss and must be prevented. Always maintain proper insulation.

\subsubsection{Management}

\begin{enumerate}[leftmargin=*]
\item Call emergency services
\item Move person to warm environment if possible; if outdoors, provide shelter from wind and moisture
\item Remove wet clothing and replace with dry blankets, clothing, or sleeping bags
\item Warm the core first (chest, neck, head, groin) using skin-to-skin contact under blankets or warm packs wrapped in cloth -- avoid direct heat on extremities as this can cause dangerous redistribution of cold blood to the core
\item Give warm sweet liquids if person is conscious and able to swallow (no alcohol or caffeine)
\item Monitor breathing; if not breathing normally, begin CPR (hypothermic patients may appear dead but can often be revived with prolonged CPR and medical care)
\end{enumerate}
\textbf{WARNING:} \textbf{CAUTION:} Do NOT rub or massage frozen limbs, give alcohol, or use direct heating (hot water bottles, heating pads) on extremities rapidly, as these can cause cardiac arrhythmias or further heat loss. Handle the person gently to prevent cardiac arrest in severe cases.

\subsection{Frostbite}

Freezing of skin and underlying tissues caused by exposure to cold temperatures, most commonly affecting fingers, toes, nose, ears, and cheeks.

\begin{description}[leftmargin=!]
  \item[Frostnip:] Reversible superficial injury with white/pale skin, tingling or numbness but no tissue damage
  \item[Superficial Frostbite:] Frozen outer skin layers with white/yellowish skin, hard/waxy appearance, possibly blistering after rewarming
  \item[Deep Frostbite:] Affects all skin layers and underlying tissue; skin appears white/blue-gray, hard/bone-like sensation; blackened eschar (dead tissue) may develop later
\end{description}

\subsubsection{Management}

\begin{enumerate}[leftmargin=*]
\item Get to a warm environment quickly
\item Gently rewarm affected area using warm (not hot) water (104-107.6 degrees Fahrenheit / 40-42 degrees Celsius) for 15-30 minutes until skin regains color and feeling returns. Alternatively use body heat (place frostbitten hands in armpits).
\item Do NOT rub or massage frozen tissue -- this causes further damage
\item Do NOT use direct heat sources (heating pads, fire, stove) -- numbness means the person cannot feel burns
\item Allow blisters to heal on their own -- do not break them (high infection risk)
\item Wrap loosely in sterile gauze or clean cloth between separated fingers/toes to keep them dry and prevent sticking together
\item Give warm fluids; avoid smoking/nicotine (constricts blood vessels)
\item Seek medical attention for deep frostbite or any case involving extensive area
\end{enumerate}
\medskip\noindent\textbf{Prevention Tip.}
Layer clothing appropriately in cold weather: moisture-wicking base layer, insulating middle layer (fleece or wool), waterproof/windproof outer layer. Cover extremities (hat, mittens instead of gloves). Stay dry and change wet clothing promptly. Take frequent breaks in warm areas during extended outdoor activities.

\section{Drowning and Water Rescue}

Drowning occurs when submersion in liquid leads to respiratory impairment. Every second counts underwater.

\subsection{Rescue Safety}

\textbf{WARNING:} \textbf{CAUTION:} First and foremost ensure your own safety. Never become another victim. Use reaching assists (pole, branch) or throwing assists (rope, flotation device) whenever possible. Only enter the water as a last resort and only if trained in water rescue.

\subsubsection{Reach Throw Row Go}

\begin{enumerate}[leftmargin=*]
  \item \textbf{REACH}: Extend arm, pole, branch, or towel to pull victim to safety
  \item \textbf{THROW}: Toss a life ring, rope, flotation device, or empty plastic jug toward the victim and instruct them to grab it
  \item \textbf{ROW}: Use a boat, paddleboard, or other flotation craft to reach the victim
  \item \textbf{GO}: Enter water only as last resort; approach the victim from behind to avoid being grabbed, secure yourself with a rope if possible, bring victim to safety on their back with airway clear, provide CPR if needed
\end{enumerate}
\subsection{After Getting Victim Out of Water}

\begin{enumerate}[leftmargin=*]
  \item Call emergency services immediately (dial 112 or 108 in India) if not already done
  \item Check responsiveness and breathing
  \item If not breathing but has a pulse, begin rescue breathing (1 breath every 5-6 seconds for an adult, 1 breath every 2-3 seconds for a child/infant); recheck the pulse every 2 minutes
  \item If no pulse, begin full CPR starting with chest compressions
  \item Even if the victim seems fine after resuscitation, seek medical evaluation as "secondary drowning" (delayed pulmonary edema) can occur hours later \cite{erc_drowning}
  \item If unconscious but breathing, place in recovery position and monitor until emergency services arrive
\end{enumerate}
\section{Marine Envenomations}

Jellyfish, stingrays, sea urchins, and other marine life can sting. Reactions range from local pain to life-threatening allergic or toxic reactions.

\subsection{Jellyfish Stings}

\begin{enumerate}[leftmargin=*]
\item Get the person out of the water carefully; keep them still to reduce venom spread
\item Rinse the affected area with \textbf{seawater} (not fresh water, which can trigger further discharge of stinging cells) to wash away invisible stinging cells
\item Remove remaining tentacles with a gloved hand, towel, or tweezers -- do NOT rub or scrape the skin
\item Immerse the area in hot water (not scalding) around 104--113\textdegree F (40--45\textdegree C) for 20--45 minutes or until pain subsides, if hot water is available; otherwise apply a cold pack for pain relief
\item Vinegar is useful only for potentially lethal species such as the box jellyfish in tropical waters -- do NOT use vinegar for bluebottle (Portuguese man-of-war) stings
\item Do NOT rinse with fresh water, apply alcohol, rub the area with sand, or scrape the skin with a razor
\item Call emergency services for stings covering a large area, on the face/genitals, or with difficulty breathing, chest pain, or severe allergic symptoms
\end{enumerate}
\quicknote{Regional differences matter: in Australian waters, box jellyfish stings require vinegar and immediate emergency care, and some jellyfish stings should NOT be treated with hot water. Follow local guidance where stings are common.}

\subsection{Stingray Stings}

Stingray wounds are puncture lacerations that can bleed heavily and retain barb fragments. The venom causes severe, rapidly escalating pain.

\begin{enumerate}[leftmargin=*]
\item Immerse the wound in hot water (104--113\textdegree F / 40--45\textdegree C) as tolerated for 30--90 minutes to inactivate the venom and relieve pain
\item Control bleeding with direct pressure
\item Do NOT attempt to remove an embedded barb in the field -- immobilize and seek emergency care
\item Seek medical care for all stingray stings: they carry a high risk of infection and retained fragments
\end{enumerate}
\subsection{Sea Urchin and Coral Injuries}

\begin{itemize}[leftmargin=*]
\item Sea urchins: soak the wound in hot water for pain, then remove visible spines with tweezers -- do NOT dig. Seek care for deeply embedded spines (especially near joints), signs of infection, or if the person was exposed to a venomous species
\item Coral scrapes: wash thoroughly with soap and water, apply antiseptic, and watch for slow-healing wounds or infection
\end{itemize}

\section{Tick Bites and Lyme Disease}

Ticks attach to skin and can transmit diseases including Lyme disease. Prompt, correct removal reduces risk.

\subsubsection{Removal}

\begin{enumerate}[leftmargin=*]
\item Use fine-tipped tweezers; grasp the tick as close to the skin surface as possible
\item Pull upward with steady, even pressure -- do not twist or jerk
\item Do NOT squeeze, crush, or burn the tick, and do NOT apply petroleum jelly or nail polish (these can cause the tick to regurgitate into the wound)
\item After removal, clean the area and your hands with soap and water or alcohol
\item Save the tick in a sealed container or photograph it for identification
\end{enumerate}
\subsubsection{When to Seek Medical Care}

\begin{itemize}[leftmargin=*]
\item A red "bull's-eye" rash (erythema migrans) or any expanding rash within 30 days
\item Flu-like symptoms after a tick bite (fever, chills, headache, muscle aches, fatigue)
\item Inability to remove the whole tick
\item Symptoms of serious tick-borne disease: facial drooping, severe headache, stiff neck, joint swelling, or heart palpitations
\end{itemize}

\textbf{WARNING:} Antibiotics are most effective when started early for Lyme disease. Do not wait for symptoms if you know a tick was attached for 24 hours or more.

\section{Altitude Sickness}

Acute mountain sickness (AMS) affects people who ascend above about 8,000 feet (2,400 m) too quickly. It can progress to life-threatening forms: high-altitude cerebral edema (HACE) and high-altitude pulmonary edema (HAPE).

\begin{itemize}[leftmargin=*]
\item \textbf{Mild AMS symptoms}: headache, nausea, fatigue, dizziness, poor sleep
\item \textbf{Red flags (HACE/HAPE)}: confusion, clumsiness, staggering, severe headache, vomiting, shortness of breath at rest, cough, chest tightness, blue lips/fingernails
\end{itemize}

\subsubsection{Management}

\begin{enumerate}[leftmargin=*]
\item Stop ascending; rest and allow time for acclimatization
\item Do NOT ascend further if symptoms are present
\item Descend immediately if symptoms worsen or any red flags appear -- descent of 1,000--3,000 feet (300--1,000 m) often relieves symptoms
\item Provide oxygen if available; pain relievers such as acetaminophen or ibuprofen may help headache
\item Call emergency services and arrange evacuation for any red-flag symptoms
\item Do NOT leave a person with severe altitude illness alone, and do NOT let them descend unaccompanied
\end{enumerate}
\section{Lightning Strike}

Lightning is the most lethal weather-related cause of emergency. Strikes cause cardiac arrest, burns, and internal injury.

\begin{itemize}[leftmargin=*]
\item \textbf{Call emergency services immediately} -- every victim is potentially salvageable with prompt CPR; death usually results from cardiac arrest, not from burns
\item Move the victim to a safe area away from ongoing lightning danger -- lightning victims can be touched safely; they do not carry an electrical charge
\item Check breathing and pulse; if in cardiac arrest, begin CPR immediately and use an AED if available
\item Treat burns with cool water and cover with clean dressings
\item Seek medical evaluation for ALL lightning victims, even those who appear fine (they may have hidden injuries)
\end{itemize}

\section{Review Questions}

\begin{enumerate}[leftmargin=*]
\item How does the management of heat exhaustion differ from heat stroke?
\item Why should a hypothermic person be handled gently and their extremities not be warmed directly?
\item A jellyfish sting: what should you rinse the area with, and what water temperature may relieve pain?
\item When should someone with a tick bite seek medical care?
\item What is the first treatment for severe altitude sickness, and when is descent mandatory?
\item A person is struck by lightning and unresponsive. What should you do first?
\end{enumerate}
\index{heat stroke}
\index{heat exhaustion}
\index{heat cramps}
\index{hypothermia}
\index{frostbite}
\index{frostnip}
\index{rewarming techniques}
\index{water rescue}
\index{drowning response}
\index{reach throw row go}
\index{secondary drowning}
\index{heat-related illness}
\index{cold injury}
\index{marine envenomation}
\index{jellyfish sting}
\index{stingray sting}
\index{tick bite}
\index{Lyme disease}
\index{altitude sickness}
\index{high-altitude pulmonary edema}
\index{lightning injury}